\chapter{Related Work}
\label{ch:related_work}

% ============================================================================
% This chapter reviews prior work in sleep stage classification
% ============================================================================

\section{Sleep Stage Classification}
\label{sec:sleep_classification}

Sleep stage classification has been extensively studied in the literature. Traditional methods rely on manual scoring by trained sleep technicians following standardized criteria \cite{TODO_ADD_RECHTSCHAFFEN_KALES}. The American Academy of Sleep Medicine (AASM) defines five sleep stages: Wake (W), Non-REM stages N1, N2, N3, and REM sleep \cite{TODO_ADD_AASM}.

% TODO: Add literature review of sleep stage classification
% - Traditional manual scoring methods
% - Inter-rater reliability studies
% - Clinical importance

\section{Machine Learning for Sleep Staging}
\label{sec:ml_sleep}

Machine learning approaches to sleep stage classification can be broadly categorized into:

\subsection{Feature-Based Methods}
\label{subsec:feature_based}

Feature-based methods extract handcrafted features from EEG signals and train classical machine learning models. Common features include:

\begin{itemize}
    \item \textbf{Time-domain features:} mean, variance, skewness, kurtosis, zero-crossing rate
    \item \textbf{Frequency-domain features:} power spectral density in delta, theta, alpha, beta, and gamma bands
    \item \textbf{Time-frequency features:} wavelet coefficients, spectral entropy
\end{itemize}

Classical models used include Support Vector Machines \cite{TODO_ADD_SVM_REF}, Random Forests \cite{TODO_ADD_RF_REF}, and XGBoost \cite{TODO_ADD_XGBOOST_REF}.

\subsection{Deep Learning Methods}
\label{subsec:deep_learning}

Recent work has explored deep learning for automatic feature extraction and classification \cite{TODO_ADD_DL_REFS}. Common architectures include:

\begin{itemize}
    \item Convolutional Neural Networks (CNNs) for spatial feature extraction
    \item Recurrent Neural Networks (RNNs) and LSTMs for temporal modeling
    \item Attention mechanisms for identifying important signal segments
\end{itemize}

While deep learning achieves state-of-the-art performance on large datasets, it requires substantial data and computational resources. Feature-based methods remain competitive, especially on smaller datasets.

% TODO: Add comparison table of recent methods

\section{Cross-Validation Strategies}
\label{sec:cross_validation}

\subsection{Leave-One-Subject-Out (LOSO)}
\label{subsec:loso}

LOSO cross-validation is critical for evaluating subject-independent performance \cite{TODO_ADD_LOSO_REF}. It ensures models generalize to unseen individuals, simulating real-world deployment scenarios. Each subject is held out once as a test set while training on all remaining subjects.

LOSO is more stringent than random K-fold cross-validation, which can leak subject-specific patterns across folds. Studies show that random splits often overestimate true performance on new subjects \cite{TODO_ADD_REF}.

\subsection{Computational Challenges}
\label{subsec:computational_challenges}

LOSO requires training $N$ models for $N$ subjects, creating significant computational overhead. For 128 subjects with a training grid of 6 configurations, this requires 768 model training runs. Efficient caching strategies are essential.

% TODO: Add references to caching/efficiency work in ML

\section{Feature Selection}
\label{sec:feature_selection}

Feature selection reduces dimensionality, improves model performance, and prevents overfitting. Common methods include:

\begin{itemize}
    \item \textbf{Filter methods:} ANOVA F-test, mutual information, correlation-based selection
    \item \textbf{Wrapper methods:} Recursive feature elimination, forward/backward selection
    \item \textbf{Embedded methods:} L1 regularization, tree-based feature importance
\end{itemize}

ANOVA F-test is 200$\times$ faster than mutual information with <1\% accuracy loss for sleep stage classification \cite{TODO_ADD_BENCHMARK_REF}.

\section{Datasets}
\label{sec:datasets}

Common public datasets for sleep stage classification include:

\begin{itemize}
    \item \textbf{BOAS:} 128 subjects, 6 EEG channels (used in this thesis)
    \item \textbf{Sleep-EDF:} 197 overnight recordings, single EEG channel
    \item \textbf{SHHS:} Sleep Heart Health Study, large-scale epidemiological study
    \item \textbf{MASS:} Montreal Archive of Sleep Studies
\end{itemize}

% TODO: Add comparison table of datasets

\section{Reproducibility and Caching}
\label{sec:reproducibility}

Reproducibility is a growing concern in machine learning research \cite{TODO_ADD_REPRODUCIBILITY_REF}. Key practices include:

\begin{itemize}
    \item Fixed random seeds for deterministic results
    \item Version control for code and configuration
    \item Cache invalidation when parameters change
    \item Comprehensive experiment logging
\end{itemize}

Fingerprint-based caching using cryptographic hashes (e.g., SHA-256) ensures experiments are recomputed when parameters change while reusing results when configurations are identical \cite{TODO_ADD_CACHING_REF}.

\section{Summary}
\label{sec:related_summary}

This chapter reviewed prior work in sleep stage classification, machine learning models, cross-validation strategies, feature selection, and reproducibility practices. Our work builds on feature-based machine learning methods with LOSO cross-validation, contributing a two-tier caching system for computational efficiency.
